\documentclass[11pt,a4paper]{article}
\usepackage[paperwidth=210mm,paperheight=430mm,margin=1.5cm]{geometry}
\usepackage{graphicx}
\usepackage{float}
\usepackage{caption}
\usepackage{subcaption}
\usepackage[hidelinks]{hyperref}

% No page numbers
\pagestyle{empty}
\makeatletter
\let\ps@plain\ps@empty
\makeatother

\setlength{\parindent}{0pt}
\setlength{\parskip}{0.5em}
\captionsetup{font=small,labelfont=bf}

\newcommand{\plotdir}{../results}

\title{Generative Models Exercise: MNIST Digit Generation}
\author{}
\date{}

\begin{document}

\maketitle

This is a quick summary of the final plots from the GAN model.

The model was trained on MNIST to generate handwritten digits from random noise. To show how the results improve during training, I saved sample outputs at different epochs using the same fixed noise.

\section*{Training Curve}

\begin{figure}[H]
    \centering
    \includegraphics[width=0.72\linewidth]{\plotdir/loss_curve.png}
    \caption{Generator and discriminator loss during training. Since this is a GAN, the losses are not as easy to read as in standard supervised learning, but they still give a rough idea of how training progressed.}
\end{figure}

\section*{Generated Samples During Training}

\begin{figure}[H]
    \centering
    \begin{subfigure}[t]{0.32\linewidth}
        \centering
        \includegraphics[width=\linewidth]{\plotdir/samples_before_training.png}
        \caption{Before training}
    \end{subfigure}
    \hfill
    \begin{subfigure}[t]{0.32\linewidth}
        \centering
        \includegraphics[width=\linewidth]{\plotdir/samples_epoch_01.png}
        \caption{Epoch 1}
    \end{subfigure}
    \hfill
    \begin{subfigure}[t]{0.32\linewidth}
        \centering
        \includegraphics[width=\linewidth]{\plotdir/samples_epoch_05.png}
        \caption{Epoch 5}
    \end{subfigure}

    \vspace{0.8em}

    \begin{subfigure}[t]{0.32\linewidth}
        \centering
        \includegraphics[width=\linewidth]{\plotdir/samples_epoch_10.png}
        \caption{Epoch 10}
    \end{subfigure}
    \hfill
    \begin{subfigure}[t]{0.32\linewidth}
        \centering
        \includegraphics[width=\linewidth]{\plotdir/samples_epoch_15.png}
        \caption{Epoch 15}
    \end{subfigure}
    \hfill
    \begin{subfigure}[t]{0.32\linewidth}
        \centering
        \includegraphics[width=\linewidth]{\plotdir/samples_epoch_20.png}
        \caption{Epoch 20}
    \end{subfigure}

    \vspace{0.8em}

    \begin{subfigure}[t]{0.32\linewidth}
        \centering
        \includegraphics[width=\linewidth]{\plotdir/samples_epoch_25.png}
        \caption{Epoch 25}
    \end{subfigure}
    \hfill
    \begin{subfigure}[t]{0.32\linewidth}
        \centering
        \includegraphics[width=\linewidth]{\plotdir/samples_epoch_30.png}
        \caption{Epoch 30}
    \end{subfigure}
    \caption{Examples of generated digits during training. At the start the output is basically noise, but over time it becomes more digit-like.}
\end{figure}

\end{document}